
\IEEEraisesectionheading{\section{Introduction}\label{sec:introduction}}

\IEEEPARstart{C}{onvolutional} neural networks (CNNs) have proven to be useful models for tackling a wide range of visual tasks \cite{krizhevsky_nips2012alex, toshev_cvpr2014pose, long_cvpr2015fcn, ren_nips2015fasterrcnn}. At each convolutional layer in the network, a collection of filters expresses neighbourhood spatial connectivity patterns along input channels---fusing spatial and channel-wise information together within local receptive fields.  By interleaving a series of convolutional layers with non-linear activation functions and downsampling operators, CNNs are able to produce image representations that capture hierarchical patterns and attain global theoretical receptive fields.
A central theme of computer vision research is the search for more powerful representations that capture only those properties of an image that are most salient for a given task, enabling improved performance.  As a widely-used family of models for vision tasks, the development of new neural network architecture designs now represents a key frontier in this search.
Recent research has shown that the representations produced by CNNs can be strengthened by integrating learning mechanisms into the network that help capture spatial correlations between features.  
One such approach, popularised by the Inception family of architectures \cite{szegedy_cvpr2015googlenet, ioffe_icml2015bn}, incorporates multi-scale processes into network modules to achieve improved performance. Further work has sought to better model spatial dependencies \cite{bell_cvpr2016ionet, newell_eccv2016hourglass} and incorporate spatial attention into the structure of the network \cite{max_nips2015stn}.

%--------------------------------------------------------------------------
\begin{figure*}[t]
\begin{center}
\includegraphics[width=0.85\textwidth]{figures/pipeline.pdf}
\end{center}
\caption{A Squeeze-and-Excitation block.}
\label{fig:splash}
\end{figure*}

In this paper, we investigate a different aspect of network design - the relationship between channels.  We introduce a new architectural unit, which we term the {\it Squeeze-and-Excitation} (SE) block, with the goal of improving the quality of representations produced by a network by explicitly modelling the interdependencies between the channels of its convolutional features. To this end, we propose a mechanism that allows the network to perform feature recalibration, through which it can learn to use global information to selectively emphasise informative features and suppress less useful ones.

The structure of the SE building block is depicted in Fig.~\ref{fig:splash}.  For any given transformation $\FF_{tr}$ mapping the input $\XX$ to the feature maps $\UU$ where $\UU\in \mathbb{R}^{H \times W \times C}$, e.g. a convolution, we can construct a corresponding SE block to perform feature recalibration. 
The features $\UU$ are first passed through a \textit{squeeze} operation, which produces a channel descriptor by aggregating feature maps across their spatial dimensions ($H \times W$). The function of this descriptor is to produce an embedding of the global distribution of channel-wise feature responses, allowing information from the global receptive field of the network to be used by all its layers.  The aggregation is followed by an \textit{excitation} operation, which takes the form of a simple self-gating mechanism that takes the embedding as input and produces a collection of per-channel modulation weights. These weights are applied to the feature maps $\UU$ to generate the output of the SE block which can be fed directly into subsequent layers of the network.

It is possible to construct an SE network (SENet) by simply stacking a collection of SE blocks. Moreover, these SE blocks can also be used as a drop-in replacement for the original block at a range of depths in the network architecture (Section~\ref{subsec:ablation-stages}). While the template for the building block is generic, the role it performs at different depths differs throughout the network.  In earlier layers, it excites informative features in a class-agnostic manner, strengthening the shared low-level representations.  In later layers, the SE blocks become increasingly specialised, and respond to different inputs in a highly class-specific manner (Section~\ref{subsection:role-of-excitation}).  As a consequence, the benefits of the feature recalibration performed by SE blocks can be accumulated through the network.

The design and development of new CNN architectures is a difficult engineering task, typically requiring the selection of many new hyperparameters and layer configurations.  By contrast, the structure of the SE block is simple and can be used directly in existing state-of-the-art architectures by replacing components with their SE counterparts, where the performance can be effectively enhanced. SE blocks are also computationally lightweight and impose only a slight increase in model complexity and computational burden. 

To provide evidence for these claims, we develop several SENets and conduct an extensive evaluation on the ImageNet dataset \cite{russakovsky_ijcv2015imagenet}. We also present results beyond ImageNet that indicate that the benefits of our approach are not restricted to a specific dataset or task.
By making use of SENets, we ranked first in the ILSVRC 2017 classification competition. Our best model ensemble achieves a $2.251\%$ top-5 error on the test set\footnote{\scriptsize{\url{http://image-net.org/challenges/LSVRC/2017/results}}}. This represents roughly a $25\%$ relative improvement when compared to the winner entry of the previous year (top-5 error of $2.991\%$).